\section{The coefficient sphere and curves through a point}
\label{sec:families}

The sphere of almost-complex structures associated with a definite triple
has one more useful feature than an individual tamed structure: its
symplectic perturbations wind once around the Seiberg--Witten wall.
The use of this winding sphere to construct an elliptic fibration
realizes the program proposed in \cite[Section~7.4.1]{Li2010SCY}.
We calculate the resulting families invariant, including the line bundle
involved in evaluating a spinor at a point.  Only the existence and
compactness direction of Taubes's theorem is needed.

Throughout this section, curve classes are identified with their
Poincare duals.  An integral class written in
\(\Lambda=H^2(X;\Z)/\mathrm{torsion}\) may be lifted to
\(H^2(X;\Z)\) when specifying a line bundle.  Curve configurations and
their total classes will only be required modulo torsion.
By \Cref{prop:topology},
\begin{equation}
 \label{eq:families-topology}
 b^+=3,\qquad
 (b_1,\sigma)=(0,-16)\ \text{or}\ (4,0),\qquad
 2\chi+3\sigma=0.
\end{equation}
Let \(J_v\), \(v\in S^2\), be the coefficient sphere of
\Cref{lem:twistor-degree}, with sign chosen so that
\(B_v=v\cdot\omega\) tames \(J_v\).  Write
\(P=\operatorname{span}\{[\omega_1],[\omega_2],[\omega_3]\}\) for the
original positive period plane.

A positive curve configuration is a finite sum of irreducible closed
pseudoholomorphic curves with positive integer multiplicities.  The
main output of the section is the following existence statement.

\begin{theorem}
 \label{thm:family-curves}
 Let \(F\in\Lambda\) be nonzero with \(F^2=0\), and write
 \(F^+=\operatorname{pr}_P F\).  For every \(p\in X\),
 there exists a finite positive \(J_v\)-holomorphic configuration
 of total class \(F\) whose support contains \(p\).
 Its parameter is necessarily the coefficient direction
 \[
  v=\frac{F^+}{|F^+|},
 \]
 where \(P\) is identified with \(\R^3\) by the normalized
 ordered period basis.  Thus the almost-complex structure is the
 same for all prescribed points.
\end{theorem}

The proof uses the wall-crossing calculation to force a zero of the
first spinor component at a prescribed point, then retains that point
in a large-parameter curve limit.  It produces configurations; the
choice of class in \Cref{sec:lattice-chamber} will make each of them
one reduced irreducible curve.

\subsection{A smooth compatible family}

\begin{lemma}
 \label{lem:compatible-family}
 There is a smooth family of closed \(J_v\)-compatible forms
 \(\eta_v\).  Orthogonal projection to the original positive plane
 \(P\) satisfies
 \[
  \operatorname{pr}_P[\eta_v]=a(v)[B_v],\qquad a(v)>0.
 \]
 Consequently the normalized period projection of this family has
 degree one in the coefficient coordinates.
\end{lemma}

\begin{proof}
Fix a metric \(g_0\) in the conformal class determined by the original
triple.  All \(J_v\) are orthogonal for \(g_0\), and the three closed
self-dual forms \(\omega_i\) span its self-dual harmonic forms because
\(b^+=3\).  Denote by \(R_v\) the projection to the
\(J_v\)-anti-invariant two-forms, and consider
\[
 L_v^{\mathrm{ell}}=R_v\dd\dd^*
 \quad\text{on the anti-invariant forms},
\]
where the adjoint is taken with respect to \(g_0\).
This is the nonnegative elliptic operator in
\Cref{lem:closed-lift}.  Its kernel consists exactly of the closed
anti-invariant forms: its quadratic form is \(\|\dd^*\xi\|_2^2\),
and an anti-invariant form is self-dual, so coclosedness implies
closedness.  Its kernel is therefore precisely
\[
 \left\{\sum_{i=1}^3 a_i\omega_i:a\cdot v=0\right\},
\]
a smooth vector bundle of rank two over \(S^2\).

Let \(H_v\) be the projection onto this explicit kernel.  In local
parameter coordinates, identify the varying anti-invariant bundles by
their orthogonal projections.  The kernel projection is then smooth,
and \(L_v^{\mathrm{ell}}+H_v\) is invertible from \(H^{s+2}\)
to \(H^s\).  The inverse identity and elliptic regularity give
smooth dependence on each Sobolev scale, as in the parameter statement
of \Cref{lem:closed-lift}.  Thus the generalized inverse is
\[
 G_v=(L_v^{\mathrm{ell}}+H_v)^{-1}(\Id-H_v),
\]
and
\begin{equation}
 \label{eq:families-closed-correction}
 D_v=\dd\dd^*G_v+H_v,
 \qquad R_vD_v=\Id,
 \qquad \dd D_v=0
\end{equation}
is a smooth family of closed lifts.

For any \(v_0\), \Cref{thm:tamed-compatible} supplies a compatible
form \(\eta\) for \(J_{v_0}\).  For \(v\) near \(v_0\), set
\(\eta_v^{(v_0)}=\eta-D_vR_v\eta\).  These forms are closed and
\(J_v\)-invariant, equal \(\eta\) at \(v_0\), and remain positive
after shrinking the parameter neighborhood.  A finite parameter
partition of unity patches these local families.  Since its
coefficients depend only on \(v\), this operation preserves
closedness on \(X\); positivity and invariance are preserved by
convexity.

The class \([\eta_v]\) is orthogonal to the classes of the two
anti-invariant forms, so its projection to \(P\) is a multiple of
\([B_v]\).  Moreover
\(\int_X\eta_v\wedge B_v>0\): the invariant part of \(B_v\) is
positive, and its anti-invariant part wedges trivially with
\(\eta_v\).  Since \([B_v]^2=1\), the multiple is positive.
\end{proof}

\subsection{The fixed spin-c structure and the equations}

Fix the canonical \(\mathrm{Spin}^c\) structure of one \(J_v\),
and denote its isomorphism type by \(\mathfrak s_0\).
The structures \(J_v\) form a connected family, so all their
canonical structures have this same fiberwise isomorphism type.
For \(E\in H^2(X;\Z)\), let
\(\mathfrak s_E=\mathfrak s_0\otimes L_E\), where
\(c_1(L_E)=E\).
We use the extension of this fixed structure
over the contractible space of metrics; there is no twist from the
parameter base.  Its determinant class and dimensions are
\begin{equation}
 \label{eq:families-dimensions}
 c_1(\det\mathfrak s_E)=2E,
 \qquad d(\mathfrak s_E)=E^2,
 \qquad \operatorname{ind}_{\C}\not D_E
            =\frac{4E^2-\sigma}{8}.
\end{equation}
Fix a homology orientation once and for all.

Given a compatible pair \((\eta,J)\), put
\(g(u,w)=\eta(u,Jw)\).  Then \(\eta\) is self-dual and has length
\(\sqrt2\).  For a fixed \(\mathrm{Spin}^c\) structure with positive
spinor bundle \(W^+\) and determinant line \(L\), we use the
following normalization of the symplectic Seiberg--Witten equations:
\begin{equation}
 \label{eq:families-taubes-equations}
 \not D_{\mathcal A}\psi=0,
 \qquad
 F_{\mathcal A}^{+}
   =r\,\mathfrak q_g(\psi)-\frac{ir}{4}\eta,
 \qquad r\ge1.
\end{equation}
Here \(\mathcal A\) is a unitary determinant connection and
\(\mathfrak q_g\) is the quadratic map with the normalization of
\cite[Equation~(2.9)]{Taubes1999}.  The spinor is normalized by
\(r^{-1/2}\) relative to the usual unscaled equations.

Clifford multiplication by \(\eta\) splits
\(W^+=\mathcal E\oplus K_J^{-1}\mathcal E\), with
\(\psi=(\alpha,\beta)\); the first summand is its
\(-i\sqrt2|\eta|=-2i\) eigenspace.  Equivalently,
\begin{equation}
 \label{eq:families-first-component}
\alpha=\frac12\left(\Id+\frac i2c^+(\eta)\right)\psi.
\end{equation}
This is the convention of \cite[Equation~(3.17) and
Section~6]{Taubes1999}, which is used in its zero-support proof.
In particular this is an intrinsic projection in a fixed
\(\mathrm{Spin}^c\) structure, and its zero set does not depend on
a local identification with the canonical structure.

For the coefficient family \(J_v\) and the fixed structure
\(\mathfrak s_E\), the Spin-c twisting action, together with the
actual homotopy of the canonical structures, identifies this first
summand on every fiber with a line of integral Chern class \(E\).
This identification uses the Spin-c structure itself: equality of
determinant Chern classes alone would lose possible two-torsion.

\begin{lemma}[Compact-family Taubes compactness with incidence]
 \label{lem:families-taubes-compactness}
 Let \((\eta_n,J_n,g_n)\) be compatible symplectic backgrounds on
 \(X\) converging in \(C^\infty\) to
 \((\eta_\infty,J_\infty,g_\infty)\), with
 \(g_n=\eta_n(\cdot,J_n\cdot)\).  Fix a
 \(\mathrm{Spin}^c\) structure.  Suppose that
 \(r_n\to\infty\) and that \((\mathcal A_n,\psi_n)\) solves
 \Cref{eq:families-taubes-equations} for the \(n\)-th background.
 If the first eigenline has Chern class \(E\) on each fiber, there
 is a subsequence and a possibly empty finite positive
 \(J_\infty\)-holomorphic configuration
 \[
  \mathcal C=\sum_a m_a C_a,\qquad m_a\in\N,
  \qquad [\mathcal C]=E\quad\text{in }H^2(X;\R).
 \]
 Every limit of points \(p_n\in\alpha_n^{-1}(0)\) belongs to
 \(\spt\mathcal C\).  In particular, if \(E=0\) over \(\R\),
 then no such sequence of zeros exists.  Over any compact smooth
 family of compatible backgrounds, all first components in this
 zero-class case are nowhere zero for sufficiently large \(r\),
 with one threshold for the whole family.
\end{lemma}

The proof, including the passage to the limiting almost-complex structure
and retention of the point constraint, is given in
\Cref{subsec:families-varying-compactness}.

\begin{remark}
 \label{rem:families-exact-equations}
Equation~\eqref{eq:families-taubes-equations} uses no additional
bounded curvature term in its perturbation.  This is the exact
normalization of \cite[Theorem~2.2]{Taubes1999}.  Adding the usual bounded
canonical-curvature term gives the same large-parameter chamber,
but such a change is unnecessary here.  Canonical determinant
connections will enter the curvature current in the compactness proof;
only their restrictions and curvatures along \(X\) occur.  No
connection in parameter directions is required by the equations.
\end{remark}

\subsection{The spherical wall coefficient}

For \(p\in X\), quotient the irreducible configuration space first
by the gauge transformations equal to one at \(p\).  Its remaining
circle bundle defines a degree-two class \(U_p\).  Choose its sign
so that, on a reducible link given by projectivizing Dirac kernels,
it restricts to \(c_1(\mathcal O(1))\).
The spinor gauge action has weight one in this convention.

Write
\[
 \Pic^0(X)=H^1(X;\R)/\operatorname{im}\bigl(H^1(X;\Z)
                                      \longrightarrow H^1(X;\R)\bigr)
\]
for the torus of topologically trivial unitary flat line bundles.
It is the identity component of the space of unitary flat line bundles;
the universal line over \(X\times\Pic^0(X)\) is normalized at \(p\).

\begin{lemma}[Critical spherical wall crossing]
 \label{lem:families-wall-crossing}
 Suppose \(E^2\ge-2\), and put \(q=(E^2+2)/2\).
 For the product manifold family over \(S^2\), with the fixed
 \(\mathfrak s_E\) just specified, the difference between two
 chamber evaluations of \(U_p^q\) is
 \begin{equation}
  \label{eq:families-wall-formula}
  \pm\deg(\text{difference of wall sections})
       \int_{\Pic^0(X)}
          s_{b_1/2}(\operatorname{Ind}\not D_E).
 \end{equation}
 Here \(s\) denotes the inverse total Chern class, and the Dirac
 family over \(\Pic^0(X)\) uses the universal flat twisting line
 normalized at \(p\).  The constant family evaluation is zero.
 For the large-parameter family
 \Cref{eq:families-taubes-equations} associated with \(\eta_v\),
 the degree difference from a constant family has absolute value
 one.
\end{lemma}

\begin{proof}
The critical wall-crossing theorem and its section-class version are
\cite[Theorem~4.10 and Proposition~4.11]{LiLiu2001}.  We check the
normalization and degree relevant here.
A transverse homotopy between two families meets the harmonic wall
at isolated parameters in \(S^2\times[0,1]\), because its normal
rank is \(b^+=3\).  Above such a parameter the reducibles form the
torus \(\Pic^0(X)\).  Represent the complex Dirac index on this
torus as \(K-O\), with ranks \(k,o\), using a stabilization when
necessary.  Write \(x=c_1(\mathcal O(1))\) on the projective
bundle of lines \(\pi:\mathbb P(K)\to\Pic^0(X)\).
The obstruction bundle contributes
\(e(\mathcal O(1)\otimes O)=\sum_i x^{o-i}c_i(O)\).
Projective pushforward therefore gives
\begin{align}
 \label{eq:families-projective-pushforward}
 \pi_*\left(x^q\sum_i x^{o-i}c_i(O)\right)
 &=\sum_i s_{q+o-(k-1)-i}(K)c_i(O)\\
 &=s_{b_1/2}(K-O),\notag
\end{align}
since \Cref{eq:families-topology,eq:families-dimensions} imply
\(q+o-(k-1)=b_1/2\).
The normal crossing signs sum to the degree difference, giving
\Cref{eq:families-wall-formula}.

The based-gauge normalization makes the universal flat line trivial
at \(p\).  Thus no Picard-base line is added to \(x\).
Writing a determinant connection as a canonical connection plus
twice a twisting connection does not alter this conclusion: the
remaining circle acts on the spinor with weight one, as used in
\Cref{eq:families-projective-pushforward}.

For a constant regular family, negative ordinary expected dimension
gives an empty moduli space.  Otherwise the moduli space is the
product of an ordinary \(E^2\)-dimensional moduli space with
\(S^2\), while \(U_p^q\) is pulled back from the former and has
degree \(E^2+2\).  Its pairing is therefore zero.

Every maximal positive harmonic plane projects isomorphically to
\(P\): the kernel would lie in the negative definite space
\(P^\perp\).  This identifies the harmonic-wall bundle with the
fixed oriented vector space \(P\).  For the perturbation in
\Cref{eq:families-taubes-equations}, its wall section has leading
term \((r/4)[\eta_v]\), with a bounded shift from \(2E\).
By \Cref{lem:compatible-family}, it is nonzero for all sufficiently
large \(r\) and its normalization has degree one.  The same is
true after sufficiently small generic family perturbations.
\end{proof}

All families pairings below are computed with such generic
perturbations in the indicated chamber.  The following elementary
compactness observation specifies how we return to the exact
equations.  At a fixed \(r\), consider a sequence of backgrounds
in a compact smooth parameter family and a sequence of smooth
perturbations converging in \(C^\infty\).  Any corresponding
sequence of solutions has, modulo gauge, a smoothly convergent
subsequence solving the limiting equations.
This follows from the spinor Weitzenbock bound, the abelian
curvature equation, gauge fixing and elliptic bootstrapping; see
\cite[Section~2.1]{LiLiu2001}.  Thus nonemptiness for perturbations
tending to zero gives a solution of the exact equations.  If all
the perturbed solutions chosen satisfy \(\alpha(p)=0\), that
condition is retained.  In using
\Cref{lem:families-taubes-compactness} we first take this
perturbation limit at each fixed \(r\), and only then let
\(r\to\infty\).

\subsection{The evaluation line over the coefficient sphere}

For \(E^2=0\), write \(\operatorname{FSW}_E(U_p)\) for the
spherical pairing of \(U_p\) in the large-parameter chamber, and
\(\operatorname{SW}_X(E)\) for the ordinary dimension-zero invariant
of \(\mathfrak s_E\).  Denote by \(\pi\) the projection of the
regular family moduli space to \(S^2\).

Let \(h\in H^2(S^2;\Z)\) be the oriented generator.  The first
eigenline in the fixed spin-c family is a line bundle
\(\mathcal E_E\to X\times S^2\), whose restriction to each
fiber has Chern class \(E\).

\begin{lemma}
 \label{lem:families-evaluation-line}
 There is an integer \(e\), with \(|e|=1\), such that
 \[
  c_1(\mathcal E_E)=E+eh.
 \]
 If \(E^2=0\), first-component evaluation at \(p\) on a regular
 two-dimensional family moduli space is a section of a line bundle
 with first Chern class \(U_p+e\pi^*h\).  Consequently, if that
 evaluation is nowhere zero,
 \begin{equation}
  \label{eq:families-evaluation-identity}
  \operatorname{FSW}_E(U_p)+e\operatorname{SW}_X(E)=0.
 \end{equation}
 If the left side is nonzero, the exact large-parameter family has
 a solution with \(\alpha(p)=0\).
\end{lemma}

\begin{proof}
Since \(H^1(S^2)=0\), the Kunneth decomposition shows that
\(c_1(\mathcal E_E)=E+eh\) for some integer \(e\).
At \(p\), deform the compatible metrics to \(g_0\) through
\(J_v\)-Hermitian metrics.  The fixed spin-c extension over metric
space identifies the positive spinor bundles along this homotopy.
For \(g_0\), the first eigenspaces give the tautological line in
the projective spinor sphere, pulled back by the coefficient sphere
map.  The latter has degree of absolute value one by
\Cref{lem:twistor-degree}, proving \(|e|=1\).

The evaluation transforms with the same gauge weight as the
spinor.  Its sign can also be checked on a reducible link: evaluation
on a kernel line \(\ell\) is a map
\(\ell\to\mathcal E_E|_p\), hence lies in
\(\operatorname{Hom}(\ell,\mathcal E_E|_p)
=\mathcal O(1)\otimes\mathcal E_E|_p\).
On the quotient it is therefore a section of the
tensor product of the based-gauge evaluation line with the
pullback of \(\mathcal E_E|_{\{p\}\times S^2}\).
Its Chern class is \(U_p+e\pi^*h\).  Pairing \(\pi^*h\)
with the family moduli cycle gives the ordinary dimension-zero
invariant \(\operatorname{SW}_X(E)\); equivalently, restrict
the family to a regular base point.  A nowhere-zero evaluation
section trivializes its line, giving
\Cref{eq:families-evaluation-identity}.

For the last assertion, suppose the exact family had no evaluation
zero.  Fixed-\(r\) compactness would give the same absence for all
sufficiently small perturbations: otherwise a sequence of their
zeros would limit to an exact zero.  Apply the preceding identity
to a regular such perturbation to obtain a contradiction.
\end{proof}

The parameter line just computed also explains the globalization
of the spinor splitting.  The varying canonical determinant over
\(X\times S^2\) has Chern class \(-2eh\), while
\(c_1(\mathcal E_E)=E+eh\).  Hence
\[
 c_1(K^{-1}\mathcal E_E^2)=-2eh+2(E+eh)=2E,
\]
as required for the fixed determinant family.  Fiberwise canonical
identifications may therefore be used on parameter charts, with
their natural transition maps, even though the canonical family
itself is not the pulled-back fixed family.  The vertical canonical
connections and their curvature two-forms are globally defined.

\subsection{The Picard-torus calculation}

For \(b_1=0\), the integral in
\Cref{eq:families-wall-formula} equals one.  For \(b_1=4\) we
must calculate rather than assume a torus cohomology ring.
Let
\(u\in H^1(X;\Z)\otimes H^1(\Pic^0(X);\Z)\)
be the mixed Chern class of the normalized universal flat line.
The families index theorem, in the form used in
\cite[Section~4.2]{LiLiu2001}, gives
\begin{equation}
 \label{eq:families-index-character}
 \operatorname{ch}(\operatorname{Ind}\not D_E)
       =\int_X e^{E+u}\widehat A(TX),
 \qquad
 \operatorname{ch}_1=\int_X\frac{Eu^2}{2},
 \quad
 \operatorname{ch}_2=\int_X\frac{u^4}{24}.
\end{equation}
In degree four on the Picard torus,
\begin{equation}
 \label{eq:families-segre-two}
 s_2=c_1^2-c_2=\frac12\operatorname{ch}_1^2
                          +\operatorname{ch}_2.
\end{equation}

The real cohomology-ring conclusion is known for symplectic
Calabi--Yau surfaces with \(b_1=4\); see
\cite[Proposition~7.8(4)]{Li2010SCY}.  We recover it here from the
marked family calculation, which also gives the Segre pairing needed
below.

\begin{proposition}
 \label{prop:cup-product}
 If \(b_1=4\), the quadruple cup product on \(H^1(X;\R)\)
 is nonzero.  The cup-product map
 \(\bigwedge^2H^1(X;\R)\to H^2(X;\R)\) is an isomorphism,
 and every spherical homology class vanishes over \(\R\).
 For every \(E\) with \(E^2=0\),
 \[
  \int_{\Pic^0(X)}s_2(\operatorname{Ind}\not D_E)\ne0.
 \]
\end{proposition}

\begin{proof}
First take \(E=0\).  Taubes's canonical nonvanishing theorem
\cite{Taubes1994} gives
\(|\operatorname{SW}_X(0)|=1\), since \(b^+>1\).
By \Cref{lem:families-taubes-compactness}, the exact symplectic
family has no first-component zeros for all sufficiently large
\(r\).  Fixed-\(r\) compactness transfers this absence to
sufficiently small regular perturbations.  Applying
\Cref{eq:families-evaluation-identity} gives
\(|\operatorname{FSW}_0(U_p)|=1\).  On the other hand,
\Cref{lem:families-wall-crossing} and
\Cref{eq:families-index-character,eq:families-segre-two} identify
this number, up to sign, with
\[
 \int_{\Pic^0(X)}\int_X\frac{u^4}{24},
\]
because \(\operatorname{ch}_1=0\) when \(E=0\).
It follows that the quadruple cup product is nonzero.

Here is an explicit calculation of the remaining assertion.  Choose
a basis \(a_1,\ldots,a_4\) of \(H^1(X;\R)\), let
\(t_1,\ldots,t_4\) be the corresponding Picard-torus basis, and
write
\[
 \delta=\int_Xa_1a_2a_3a_4\ne0,
 \qquad u=\sum_i a_it_i.
\]
The wedge pairing on \(\bigwedge^2H^1(X;\R)\) induced by this
nonzero volume element is nondegenerate.  Thus the cup-product map
into \(H^2(X;\R)\) is injective.  Both spaces have dimension
six, by \Cref{prop:topology}, so it is an isomorphism.
We may consequently write
\(E=\sum_{i<j}e_{ij}a_ia_j\).  The graded-product signs give
\begin{align}
 \label{eq:families-exterior-calculation}
 \operatorname{ch}_2&=\delta\,t_1t_2t_3t_4,\\
 \operatorname{ch}_1&=-\delta\bigl(
 e_{34}t_1t_2-e_{24}t_1t_3+e_{23}t_1t_4\notag\\
 &\hspace{40mm}
 +e_{14}t_2t_3-e_{13}t_2t_4+e_{12}t_3t_4\bigr),\notag\\
 E^2&=2\delta(e_{12}e_{34}-e_{13}e_{24}+e_{14}e_{23}),\notag\\
 \operatorname{ch}_1^2&=\delta E^2\,t_1t_2t_3t_4.\notag
\end{align}
Thus \(E^2=0\) implies \(\operatorname{ch}_1^2=0\), and
the Segre integral equals the same nonzero integral of
\(\operatorname{ch}_2\) found for \(E=0\).

Finally, a map \(S^2\to X\) pulls back every class in
\(H^1(X;\R)\) to zero, and therefore pulls back every product
of two such classes to zero.  These products span \(H^2(X;\R)\),
so its homology class is zero by Poincare duality.
\end{proof}

\subsection{Curves through a point and exclusion of roots}

\begin{proof}[Proof of \Cref{thm:family-curves}]
The projection \(F^+\) is nonzero: otherwise \(F\) would be a
nonzero square-zero vector in the negative definite space
\(P^\perp\).  Lift \(F\) to an integral class and use the
fixed spin-c structure \(\mathfrak s_F\).

Its ordinary invariant is zero.  Indeed, choose the opposite
coefficient direction \(w=-F^+/|F^+|\), so that
\(F[B_w]<0\).  If \(\operatorname{SW}_X(F)\ne0\),
ordinary invariance and fixed-parameter monopole compactness would
give solutions of \Cref{eq:families-taubes-equations} for arbitrarily
large \(r\) at this fixed background.  Taubes compactness would
give a nonempty positive \(J_w\)-holomorphic configuration of
class \(F\), contradicting its negative \(B_w\)-area.

The spherical \(U_p\) pairing is nonzero.  For \(b_1=0\),
this follows directly from \Cref{lem:families-wall-crossing};
for \(b_1=4\), use \Cref{prop:cup-product} as well.  Since the
ordinary invariant vanishes, \Cref{lem:families-evaluation-line}
forces a first-component zero over \(p\) in the exact family for
each sufficiently large \(r\).  Extracting a convergent parameter
subsequence and applying
\Cref{lem:families-taubes-compactness} gives a configuration through
\(p\) in some \(J_v\), with total class \(F\).

The two closed forms whose coefficients are perpendicular to \(v\)
are \(J_v\)-anti-invariant and restrict to zero on every component
of this configuration.  Thus \(F^+\) is parallel to \([B_v]\).
Its \(B_v\)-area is positive, so the proportionality factor is
positive, which proves the asserted direction.
\end{proof}

The unmarked wall calculation also excludes the square-minus-two classes
needed in the lattice argument.  This is a separate consequence of
\Cref{lem:families-wall-crossing}.

\begin{proposition}
 \label{prop:root-exclusion}
 If \(b_1=0\), no class \(E\in\Lambda\) with \(E^2=-2\)
 is orthogonal to \(P\).
\end{proposition}

\begin{proof}
Lift \(E\) to an integral line-bundle class.  Its ordinary expected
dimension is \(-2\), and its spherical family has dimension zero.
By \Cref{lem:families-wall-crossing}, its large-parameter unmarked
family invariant is \(\pm1\).  For each sufficiently large
\(r\), take small regular perturbations and then a fixed-\(r\)
limit to obtain an actual solution of
\Cref{eq:families-taubes-equations} at some parameter \(v_r\).
Pass to a sequence with \(v_r\to v\), and apply
\Cref{lem:families-taubes-compactness}.  The limit is a nonempty
positive \(J_v\)-holomorphic configuration of class \(E\).
Every nonconstant component has positive \(B_v\)-area, since
\(B_v\) tames \(J_v\).  This contradicts
\(E[B_v]=0\), which follows from \(E\perp P\).
\end{proof}

\begin{remark}
 \label{rem:families-period-stability}
The results of this section can be reapplied to any other definite
closed triple with the same ordered periods.  The positive plane
\(P\), the excluded roots, and the prescribed direction of \(F\)
then remain the same.  In particular, the curve-through-every-point
conclusion remains available after the small exact perturbations
used below.  At this stage the curves may have several components
or multiplicities; their reduction and irreducibility will follow
from the choice of \(F\) in the next section.
\end{remark}

\subsection{Compactness under varying backgrounds}
\label{subsec:families-varying-compactness}

We now prove the compact-family input used above.  The uniform
Seiberg--Witten estimates produce a limiting curvature current and
retain the zero sets.  To recognize this limit as curves for
\(J_\infty\), the remaining task is to construct an integer
intersection assignment and prove its positivity on
\(J_\infty\)-holomorphic disks.  This is where variation of the
background must be addressed explicitly.

\begin{proof}[Proof of \Cref{lem:families-taubes-compactness}]
For a fixed background, this is the compactness and incidence part of
Taubes's existence theorem.  For the equations in
\Cref{eq:families-taubes-equations}, use
\cite[Theorem~2.2]{Taubes1999}, with its zero-support conclusion
as proved in Lemma~7.1 and Equation~(7.5) for the first component
defined in Equation~(3.17) of that paper.  The symplectic form has
empty zero locus, so the transverse-zero hypothesis is vacuous.
The class formula is
\[
 2[\mathcal C]=c_1(L\otimes K_J).
\]
The splitting above gives \(L=K_J^{-1}\mathcal E^2\), so this is
exactly \([\mathcal C]=E\) over \(\R\).
The original fixed-background symplectic statement also includes
incidence with prescribed closed sets
\cite[Theorem~1.3]{Taubes1996}; the curve-recognition argument is
understood in its corrected form \cite{Taubes2000}.

\smallskip
\noindent\emph{Uniform estimates and support convergence.}
We explain the uniformity needed here.  Smooth convergence provides a
common positive normal-coordinate radius, uniform ellipticity, and
uniform bounds on all derivatives of the metrics, forms, and induced
canonical connections.  It also bounds the topological terms
\(c_1(L)[\eta_n]\).  In the following estimates, constants are
independent of \(n\) and \(r\) after these quantities have been
bounded.  For readability write \(w=1-|\alpha|^2\) and
\(w_+=\max(w,0)\), and use the projected spinor covariant
derivatives.  The estimates used in the
compactness proof include
\begin{align}
 \label{eq:families-taubes-basic}
 |\psi|^2&\le1+Cr^{-1},
 & |\beta|^2&\le Cr^{-1}w_++Cr^{-2},\\
 \label{eq:families-taubes-curvature}
 |F_{\mathcal A}^+|&\le\frac{r}{2\sqrt2}w_++C,
 & |F_{\mathcal A}^-|&\le
   \frac{r}{2\sqrt2}(1+Cr^{-1/2})w_++C,\\
 \label{eq:families-taubes-gradient}
 |\nabla_{\mathcal A}\alpha|^2
       +r|\nabla_{\mathcal A}\beta|^2&\le Crw_++C,
 & r\int_X|1-|\psi|^2|\,\vol_g&\le C.
\end{align}
These are the specialization to a nowhere-vanishing symplectic form of
\cite[Lemma~3.1, Proposition~3.1, Equation~(3.27), and
Propositions~3.2--3.3]{Taubes1999}.
The global pointwise bound follows directly from the spinor
Weitzenbock formula and the maximum principle.  The remaining bounds
use its two eigenline projections, the curvature equation and its
derivative, and the scalar Green kernel.  All the geometric
coefficients in those arguments are uniformly bounded on the chosen
charts.  In particular,
\(\int w_+\le\int|1-|\psi|^2|+\int|\beta|^2\).
Integrating the bound for \(\beta\) and absorbing
\(Cr^{-1}\int w_+\) for large \(r\) therefore bounds
\(r\int w_+\).  The inequality \(w\ge-Cr^{-1}\), followed by
\Cref{eq:families-taubes-curvature}, then gives
\begin{equation}
 \label{eq:families-taubes-mass}
 r\int_X|w|\,\vol_g+\int_X|F_{\mathcal A}|\,\vol_g\le C.
\end{equation}

For the nonnegative energy
\(\mathscr E_r(U)=r\int_U|1-|\psi|^2|\,\vol_g\), the
monotonicity estimate gives uniform constants \(c,C,s_0>0\) such
that, whenever \(Cr^{-1/2}\le s\le s_0\),
\begin{equation}
 \label{eq:families-taubes-monotonicity}
 \mathscr E_r(B_s(x))\le Cs^2,
 \qquad
 x\in\alpha^{-1}(0)\ \Longrightarrow\quad
 \mathscr E_r(B_s(x))\ge cs^2.
\end{equation}
This is \cite[Proposition~4.1]{Taubes1999}, with its universal
normalizing factor absorbed into the constants.  Consequently the
zero sets have covers by at most \(Cs^{-2}\) such balls.

Here is also the scale and derivative dependence in the local
compactness step.  Use normal coordinates for each individual
\(g_n\), rescale by \(r_n^{1/2}\), and compare on a ball of fixed
radius \(R\).  After identifying the center with Euclidean
\(\C^2\), the rescaled coefficients satisfy
\begin{align}
 \label{eq:families-taubes-rescaling}
 \|g_n'-g_{\mathrm E}\|_{C^0(B_R)}&\le C_Rr_n^{-1},
 &\|\partial^kg_n'\|_{C^0(B_R)}&\le C_{R,k}r_n^{-k/2}
       &&(k\ge2),\\
 \|\eta_n'-\eta_n'(0)\|_{C^0(B_R)}&\le C_Rr_n^{-1/2},
 &\|\partial^k\eta_n'\|_{C^0(B_R)}&\le C_{R,k}r_n^{-k/2}
       &&(k\ge1).\notag
\end{align}
The first derivatives of \(g_n'\) are \(O_R(r_n^{-1})\).
The scale in \cite[Equation~(5.1)]{Taubes1999} is
\((r_n|\eta_n|)^{1/2}=2^{1/4}\sqrt{r_n}\).
A fixed dilation therefore converts these coordinates to that
paper's flat-model normalization; all estimates below absorb this
fixed factor, and curvature integrals are unchanged by it.
These estimates follow by Taylor expansion in each metric's own
normal coordinates; compare \cite[Section~6, Step~4]{Taubes1999}.
For every fixed \(R,k\), the rescaled equations therefore have
uniform coefficient bounds through order \(k\).  Gauge fixing and
interior elliptic estimates give the rescaled local approximation in
any prescribed \(C^k\) norm by the flat model, as in
\cite[Proposition~5.2]{Taubes1999}.  Its quantifiers are: for each
\(R,k,\eps>0\), a common lower bound on \(r\) makes the
approximation error smaller than \(\eps\).  We retain smooth
background convergence, so every derivative order required in this
argument is available; no assertion about a minimal finite
regularity threshold is needed.

Let \(\mathcal A_{\mathrm{can},n}\) be the canonical connection
on \(K_{J_n}^{-1}\).  The first eigenline connection \(a_n\)
has curvature
\(F_{a_n}=\tfrac12(F_{\mathcal A_n}
 -F_{\mathcal A_{\mathrm{can},n}})\).
The corresponding closed currents
\begin{equation}
 \label{eq:families-taubes-current}
 T_n(\zeta)=\frac{i}{2\pi}\int_XF_{a_n}\wedge\zeta
\end{equation}
have uniformly bounded mass by
\Cref{eq:families-taubes-mass}.  Their periods on closed test forms
are those of \(E\).
On the complement of the first-component zero set, its unit section
identifies the determinant with the canonical determinant.  In that
identification the further estimate is
\begin{equation}
 \label{eq:families-taubes-off-zeros}
 |\mathcal A-\mathcal A_{\mathrm{can}}|
   +|F_{\mathcal A}-F_{\mathcal A_{\mathrm{can}}}|
 \le Cr^{-1}+Cr\exp\bigl(-\sqrt r\,
          \dist(x,\alpha^{-1}(0))/C\bigr)
\end{equation}
when the distance is at least \(r^{-1/2}\)
\cite[Proposition~6.1]{Taubes1999}.

The ball covers, the lower bound at zeros, and
\Cref{eq:families-taubes-off-zeros} now give a subsequence whose
zero sets converge to the support \(Z_*\) of a weak limit of
\(T_n\); the support has finite two-dimensional Hausdorff measure.
These are the arguments of
\cite[Equations~(7.1)--(7.5) and Lemma~7.1]{Taubes1999}, using
exactly the uniform estimates just listed.

We give the disk argument needed when the backgrounds vary.  Write
\[
 f_n=\frac{i}{2\pi}F_{a_n}
     =\frac{i}{4\pi}(F_{\mathcal A_n}
                          -F_{\mathcal A_{\mathrm{can},n}}).
\]
The sharper estimate
\cite[Proposition~4.2, Equation~(4.9)]{Taubes1999}, specialized to
\(|\eta_n|=\sqrt2\), is
\begin{equation}
 \label{eq:families-taubes-sharp}
 |F_{\mathcal A_n}^-|
 \le \frac{r_n}{2\sqrt2}(1-|\alpha_n|^2)+C.
\end{equation}
Its constants have the same uniform geometric dependence as the
estimates above.  If \(\ell\) is a unit, positively oriented
\(J_n\)-complex bivector, then
\(\ell^+=\eta_n/2\), \(|\ell^-|=1/\sqrt2\), and the
curvature equation gives
\[
 iF_{\mathcal A_n}^+(\ell)
       =\frac{r_n}{4}(1-|\alpha_n|^2+|\beta_n|^2),
 \qquad
 iF_{\mathcal A_n}^-(\ell)
       \ge-\frac{r_n}{4}(1-|\alpha_n|^2)-C.
\]
The leading terms cancel.  The canonical curvatures are uniformly
bounded, so there is one constant \(C_0\ge1\) such that every
\(J_n\)-holomorphic disk \(u\) satisfies
\begin{equation}
 \label{eq:families-disk-lower-bound}
 u^*f_n\ge-C_0\,\dd A_u.
\end{equation}
Here \(\dd A_u\) is the induced area density, counted on the
source; the inequality also holds at critical points.  This derives
the disk bound directly for our normalization, including the factor
one half in the twisting connection.

\smallskip
\noindent\emph{The relative integer assignment.}
First construct the integer assignment independently of positivity.
A smooth disk \(u\) is admissible when its boundary misses \(Z_*\).
For large \(n\), \(u^*\alpha_n\) is nonzero on a boundary
collar and defines a relative Chern number \(k_n(u)\in\Z\).
The unit first component trivializes the eigenline there, and
\Cref{eq:families-taubes-off-zeros} gives
\begin{equation}
 \label{eq:families-relative-degree}
 \int_Du^*f_n=k_n(u)+o(1).
\end{equation}
The error is uniform for disks with bounded first derivatives whose
boundary collars have a common positive distance from \(Z_*\).
For each sufficiently large \(n\), relative Chern numbers agree
along any fixed admissible homotopy.  The same conclusion holds for
the varying homotopies below: pointwise short geodesics between
uniformly close boundary maps stay in one fixed zero-free
neighborhood of the original boundary.  Thus their pulled-back
unit sections extend over the boundary homotopy, uniformly in
large \(n\).

For an immersed disk, take small normal translates, with a smooth
averaging function in the two normal parameters.  The normal bundle
is taken over the source, so self-intersections cause no difficulty.
All these disks have the same \(k_n\).  Cut off the source near its
boundary collar; the removed flux tends to zero by
\Cref{eq:families-taubes-off-zeros}.  The averaged remaining flux is
the pairing of \(T_n\) with a fixed smooth two-form: the normal
coordinate map is a local diffeomorphism, and a finite partition of
unity pushes the compactly supported averaging forms to \(X\).
Weak convergence of \(T_n\) proves convergence of these averages.
Thus \(k_n\), being integer valued, is eventually constant.  A
general admissible smooth disk has an arbitrarily close immersed
perturbation, connected to it by an admissible homotopy, and hence
has the same conclusion.  Denote the eventual integer by \(I(u)\).
It is zero for disks missing \(Z_*\), invariant under admissible
homotopy, additive under subdivision, and multiplied by the degree
under proper disk maps; these are the corresponding properties of
relative first Chern numbers.  This is the Stokes and averaging
construction of \cite[Proposition~5.6 and Lemma~6.2]{Taubes1996}.
Only positivity for the limiting complex structure remains.

\smallskip
\noindent\emph{Disk deformations with prescribed incidences.}
We need the following elementary disk deformation fact.  For a
holomorphic map on a slightly larger disk, with no boundary
condition, the linearized Cauchy--Riemann operator has a bounded
right inverse even after prescribing values and complex-linear
first derivatives at any finite set of distinct interior source
points.  Work, for example, in \(H^k\to H^{k-1}\), \(k\ge4\),
so these evaluations are continuous.  Extend the bundle and
operator to a closed surface, and extend the target sections by a
bounded Sobolev extension.  The closed-surface operator augmented
by the finite evaluations is Fredholm.  Finitely many smooth
forcing terms supported outside the larger disk make it onto.
Indeed, an annihilator of all such exterior terms vanishes there.
Away from the specified points it is smooth by elliptic regularity
and vanishes everywhere by unique continuation for the adjoint
real Cauchy--Riemann operator
\cite[Theorem~2.78 and Proposition~3.12]{Wendl2014Lectures}.
It is therefore a sum of distributions supported at the specified
points.  The invertible principal symbol eliminates derivatives of
point masses of positive order: their highest derivative after
applying the adjoint has order at least two, whereas the evaluation
constraints have order at most one.  The remaining term at each
point is \(q\delta\).  Testing independently the complex-linear
and anti-complex-linear first derivatives eliminates both \(q\)
and the first-derivative constraint covector; testing the value
then eliminates its covector.  The annihilator is zero, proving
the assertion about exterior forcing.  A bounded right inverse of
this surjective Fredholm operator, restricted to the disk, proves
the claimed right inverse.

The implicit function theorem and interior regularity consequently
give nearby holomorphic disks for nearby almost-complex structures,
with those finite values and complex-linear derivatives prescribed.
All domain restrictions are chosen first, and all constraints lie
in the smaller disk.  The right-inverse bounds persist for nearby
operators.  Constants may depend on the chosen finite set of
points and on a fixed neighborhood of the original map; this is
harmless because these choices precede \(n\to\infty\).  In
particular, while retaining finitely many incident values, one may
choose the tangent line at each of them independently in a nonempty
open cap.  At a critical point one first chooses a sufficiently
small nonzero complex-linear derivative.  This gives a fixed small
\(J_\infty\)-holomorphic perturbation of the original disk;
its \(J_n\)-holomorphic transfers converge smoothly to that
perturbation on the smaller closed disk.  They need not converge
to the unperturbed disk, and admissible boundary homotopy is what
identifies their integers.

\smallskip
\noindent\emph{Positive flux at a retained incidence.}
Here is the positive contribution at each retained incidence.
Choose actual zeros \(p_{j,n}\) tending to the finitely many
specified values in \(Z_*\).  After a joint subsequence, rescaling
at these zeros gives centered flat models on all fixed balls.
Their first components vanish at the origin.  The quadratic energy
bound excludes the identically zero first component.  Each limiting
model has a polynomial divisor by
\cite[Proposition~5.1]{Taubes1999}.
For any nonempty open cap of complex lines through the origin,
choose a direction where the highest homogeneous part of that
polynomial is nonzero; only finitely many directions are excluded.
Along the chosen line, distance from the divisor grows linearly
at infinity.  The exponential off-divisor estimate in
\cite[Proposition~5.1]{Taubes1999} makes the boundary connection
term tend to zero.  Stokes' theorem therefore identifies the total
line flux \((i/2\pi)\int F_{a_0}\) with the positive degree
\(m_j\ge1\) of the restricted divisor, which contains the
origin.  The determinant connection in
\cite[Proposition~5.2]{Taubes1999} converges to \(2a_0\), so
this is precisely the limiting normalization of \(f_n\), with
the bounded canonical curvature disappearing under rescaling.
Choose a fixed
large radius \(R_j\) so that the boundary is zero-free and the
flux differs from \(m_j\) by less than \(1/8\).
These are the polynomial-model and line-selection steps in
\cite[Proposition~5.6, Steps~3--4]{Taubes1996}, using the flat
limits of \cite[Proposition~5.2]{Taubes1999}.

Prescribe those good tangent directions in the disk deformation
just constructed.  Choose a fixed approximation tolerance making
the flux error on each model disk less than \(1/8\); only then
take \(n\) large.  Because the prescribed derivatives are
nonzero and the disk maps have uniform higher derivative bounds,
the rescaled maps on microscopic source disks converge to the
chosen complex-linear maps.  Each resulting source core therefore
contributes more than \(m_j-1/4\ge3/4\) to the flux.
The finitely many core radii tend to zero, so the cores are
source-disjoint.  Distinct source incidences may have the same
target value: the flux is integrated on the source, and this does
not affect the argument.  The nonzero derivative sizes, the caps,
the models, \(R_j\), and the approximation tolerances are all
fixed before the final lower bound on \(n\) is imposed.

\smallskip
\noindent\emph{Positivity for the limiting complex structure.}
We now prove positivity for any admissible
\(J_\infty\)-holomorphic disk \(u\) meeting \(Z_*\),
allowing critical points and self-intersections.  Shrink its domain
inside a slightly larger disk so that every incidence lies in its
interior and its boundary collar has a fixed positive gap from
\(Z_*\).  First fix a small \(C^1\) neighborhood in which
this gap persists and all disk areas are bounded by a number
\(A_*\).  Put \(S=u^{-1}(Z_*)\).

If \(S\) is finite, retain all its points.  Choose disjoint
fixed neighborhoods of them whose total induced area, for all
sufficiently close disks, is less than \(1/(4C_0)\).  On the
compact complement the original disk has a positive distance from
\(Z_*\); shrink the allowed disk neighborhood to preserve half
that distance.  Choose the good tangent caps and the fixed small
holomorphic perturbation within this neighborhood, and transfer it
to \(J_n\), retaining the chosen zeros.  The cores contribute
at least \(3/4\) each.  Their complement inside the selected
neighborhoods contributes more than \(-1/4\) by
\Cref{eq:families-disk-lower-bound}.  On the remaining compact
domain the flux tends to zero by
\Cref{eq:families-taubes-off-zeros}.  Thus the limiting integer
in \Cref{eq:families-relative-degree} is positive.

If \(S\) is infinite, choose a finite number \(N\) of distinct
source points in it with
\(3N/4>C_0A_*+1\).  The area bound was fixed before \(N\).
Apply the same finite-constraint construction at these points,
keeping the perturbed disks in the chosen area neighborhood.
The \(N\) source-disjoint cores contribute more than \(3N/4\),
whereas their full complement contributes at least \(-C_0A_*\).
The flux is therefore bounded below by a positive number.  Its
limit is the unchanged integer \(I(u)\), so again \(I(u)>0\).
Every choice was finite and fixed before the last passage to large
\(n\); no rate relating \(J_n\to J_\infty\) and
\(r_n\to\infty\) has been imposed.

\smallskip
\noindent\emph{Curve recognition and retention of the marked point.}
In particular the positivity axiom for embedded holomorphic test
disks holds.  The resulting positive cohomology assignment permits
the corrected recognition theorem to be applied to the fixed
structure \(J_\infty\), as in
\cite[Lemmas~7.2--7.3 and Proposition~7.1]{Taubes1999}
and \cite{Taubes2000}.
The recognized normalization map is proper.  Since \(X\) is
compact, its source is compact and has finitely many components;
its locally finite exceptional set is finite as well.
The multiplicity of a resulting curve at a smooth point is the
integer \(I\) of a small transverse disk there.  The averaging
construction identifies this same integer with the transverse
coefficient of the weak curvature current: on each smooth branch,
a closed current of order zero supported there is a constant
multiple of its oriented integration current.  The remaining
difference is supported on the finite singular set, where a
closed two-current of order zero must vanish.  Thus the recognized
configuration has current equal to the weak limit of \(T_n\),
as in \cite[Proposition~7.1 and Equation~(7.12) in the proof of
Lemma~7.4]{Taubes1999}; in particular its class is
\(E\).

The support convergence retains every limit of first-component
zeros.  If \(E=0\), a retained zero would give a nonempty positive
configuration with
\(\int_{\mathcal C}\eta_\infty=E[\eta_\infty]=0\), which
is impossible.  Failure of the final uniform assertion would supply
a sequence with \(r_n\to\infty\); compactness of the parameter
space would reduce it to the situation just proved.
\end{proof}
